% New English translation from checked cote 25, site 80 (PDF page 81).
% AI-translated by GPT-6 Astra (OpenAI, ChatGPT Pro mode).
\scanpage{25/80}
\mspage{\hand{75}}

\var{\margin{|}}{\placedmargin{|}}{marginal note placed}{ega25_25_080-v0001}m, noting that, up to a change of base field, $\underline{X}^{(m)}_{\hand{\unsure{\eta|\xi}}}$ is obtained by
taking an L' of codimension m-1 through x and an H of \unsure{codimension|codimension}
1 through x, L' and H being independently generic with these properties
(i.e. in orthodox terms, we work at the generic point of the \var{\xout{\unsure{the}}}{}{typewriter strike-out}{ega25_25_080-v0003}
scheme of pairs (L',H) \var{\xout{\unsure{indeed}}}{}{typewriter strike-out}{ega25_25_080-v0004}) and taking the linear section
of X by $L'\hand{\cap} H$: we may begin by taking the section by H, which
is smooth by what precedes, and then by L', which is smooth by the induction
hypothesis. This kind of argument, already used to generalize
2.6, for example, to linear sections of arbitrary codimension m, deserves
to be made explicit once and for all in general, so that we can refer to it without
going \var{\xout{into}}{}{typewriter strike-out}{ega25_25_080-v0007} each time into the rather cumbersome details of a formal
presentation.

It remains to prove the announced corollary of 8.8 in the case m=1.
When \var{\xout{dim}}{}{typewriter strike-out}{ega25_25_080-v0008} codim T $\hand{\geqslant}$ 2, since on the other hand the hyperplane \var{\inserted{Q}}{Q}{insertion}{ega25_25_080-v0009} of $\overset{\hand{\vee}}{P}$ consisting of the
$H_{\hand{\xi}}$ such that $x\hand{\in} H_{\hand{\xi}}$ has codimension 1, its generic point \var{\inserted{$\eta$}}{$\eta$}{insertion}{ega25_25_080-v0010} cannot
belong to T, and we win. In the case codim T = 1, since T is
irreducible, we can have $\hand{\eta}\hand{\in} T$ only if Q=T,\var{\xout{\unsure{(set-theoretically)}}}{}{typewriter strike-out}{ega25_25_080-v0011}
i.e. in geometric terms \var{\xout{every}}{}{typewriter strike-out}{ega25_25_080-v0012} \var{\inserted[typed]{(assuming k alg closed, which we may do)}}{(assuming k alg closed, which we may do)}{insertion}{ega25_25_080-v0013} for every $z\hand{\in} X(k)$\var{\xout{its}}{ }{typewriter strike-out}{ega25_25_080-v0014}the tangent space
of X at z (or rather its image under $f'_{z}$) passes through x. Let us show that
this cannot occur when $\underline{X}^{\mathrm{sing}}\hand{\longrightarrow}T$ is gen étale, i.e.
when we are under the conditions of 8.8, unless f(X) = \hand{\{}x\hand{\}} and hence X
has dimension 0. Indeed, 8.7 c) (which essentially expresses the
symmetry \var{\xout{between}}{}{typewriter strike-out}{ega25_25_080-v0016} in the relation between X and its \hand{“}dual\hand{”} \var{\inserted{T}}{T}{insertion}{ega25_25_080-v0017}\hand{)} then implies that
for almost every point $z\hand{\in} X(k)$, f(z) is orthogonal to the tangent space
of T at a certain point,\var{\xout{\unsure{the|of}}}{}{typewriter strike-out}{ega25_25_080-v0018} hence (since T=Q) \var{\inserted{orthogonal}}{orthogonal}{insertion}{ega25_25_080-v0019} to Q, whence f(z)=x, whence
f(X) =\hand{\{}x\hand{\}}. This completes the proof of our second thoughts, hence of 13.2.
\tnote{25-080-1}{The French letter deletions in propriétés, utilisé and passe have no separate English effect. The doubtful canceled \unsure{l} in la is retained here; it gives no independent English correction to the. The codimension/codimemsion doubt remains marked on codimension. The source's symbol X with superscript sing, and its final reference 13.2, are retained.}

\var{\margin{80}}{\placedmargin{80}}{marginal note placed}{ega25_25_080-v0020}
